\chapter{Six-Vertex Model and the XXZ Chain}
\label{chap:six-vertex-xxz}

The two-dimensional Ising model was solvable because its transfer matrix could be
rewritten in terms of free fermions.  The six-vertex model is solvable for a
different reason.  It is generally not a free-fermion model.  Its solvability
comes from the existence of an $R$-matrix satisfying the Yang--Baxter equation.
This equation implies a commuting family of transfer matrices, and the
logarithmic derivative of this family produces the one-dimensional XXZ quantum
spin chain.

This chapter has three goals:
\begin{enumerate}[label=(\roman*)]
    \item define the six allowed local vertices and the row-to-row transfer
    matrix;
    \item introduce the six-vertex $R$-matrix and explain how the Yang--Baxter
    equation gives commuting transfer matrices; and
    \item take the Hamiltonian limit and obtain the XXZ chain with anisotropy
    parameter $\Delta$.
\end{enumerate}
The coordinate and algebraic Bethe-ansatz diagonalizations are postponed to the
following chapters.  Here the main point is the exact map
\[
    \text{six-vertex transfer matrix}
    \quad\Longrightarrow\quad
    \text{commuting transfer matrices}
    \quad\Longrightarrow\quad
    \text{XXZ Hamiltonian}.
\]

\section{Arrows, the ice rule, and six local weights}
\label{sec:six-vertex-weights-ice-rule}

The six-vertex model is a classical statistical model on a square lattice.  The
degrees of freedom are arrows living on the edges.  At each vertex four arrows
meet.  The allowed configurations obey the \emph{ice rule}: exactly two arrows
point into the vertex and exactly two arrows point out of the vertex.

Equivalently, each local vertex conserves a two-state charge.  We identify an
arrow state with a spin-$1/2$ basis vector
\begin{equation}
    \ket{\uparrow}=\begin{pmatrix}1\\0\end{pmatrix},
    \qquad
    \ket{\downarrow}=\begin{pmatrix}0\\1\end{pmatrix},
    \qquad
    V\simeq \CC^2 .
\end{equation}
A local vertex weight is then an operator on $V\otimes V$.  In the ordered basis
\begin{equation}
    \ket{\uparrow\uparrow},
    \quad
    \ket{\uparrow\downarrow},
    \quad
    \ket{\downarrow\uparrow},
    \quad
    \ket{\downarrow\downarrow},
\end{equation}
the most common six-vertex weight matrix has the block form
\begin{equation}
\label{eq:six-vertex-constant-R-matrix}
    R=
    \begin{pmatrix}
        a&0&0&0\\
        0&b&c&0\\
        0&c&b&0\\
        0&0&0&a
    \end{pmatrix}.
\end{equation}
The six nonzero matrix elements are
\begin{equation}
\label{eq:six-vertex-six-nonzero-weights}
\begin{aligned}
    &\ket{\uparrow\uparrow}\longrightarrow\ket{\uparrow\uparrow}
    &&\text{with weight }a,\\
    &\ket{\downarrow\downarrow}\longrightarrow\ket{\downarrow\downarrow}
    &&\text{with weight }a,\\
    &\ket{\uparrow\downarrow}\longrightarrow\ket{\uparrow\downarrow}
    &&\text{with weight }b,\\
    &\ket{\downarrow\uparrow}\longrightarrow\ket{\downarrow\uparrow}
    &&\text{with weight }b,\\
    &\ket{\uparrow\downarrow}\longrightarrow\ket{\downarrow\uparrow}
    &&\text{with weight }c,\\
    &\ket{\downarrow\uparrow}\longrightarrow\ket{\uparrow\downarrow}
    &&\text{with weight }c.
\end{aligned}
\end{equation}
All other matrix elements vanish.  Thus $R$ conserves the total $z$-spin on the
two lines:
\begin{equation}
\label{eq:six-vertex-U1-conservation-local}
    \left[ R,\sigma^z\otimes\id+\id\otimes\sigma^z\right]=0 .
\end{equation}
This is the algebraic form of the ice rule.

\begin{definition}[Six-vertex anisotropy parameter]
For nonzero $a$ and $b$, the standard six-vertex anisotropy parameter is
\begin{equation}
\label{eq:six-vertex-Delta-from-abc}
    \Delta
    =
    \frac{a^2+b^2-c^2}{2ab} .
\end{equation}
Multiplying all weights by a common scalar changes the normalization of the
partition function but not $\Delta$.
\end{definition}

The parameter $\Delta$ is the same anisotropy parameter that appears in the XXZ
spin chain obtained from the Hamiltonian limit.  The local weights $(a,b,c)$
control a two-dimensional classical model, whereas $\Delta$ controls the relative
strength of the $z$ interaction in the associated one-dimensional quantum chain.

\begin{remark}[Why there are six vertices]
Without the ice rule there would be $2^4=16$ possible arrow configurations at a
four-valent vertex.  The constraint ``two in and two out'' leaves six.  In the
spin notation this is exactly the statement that the local weight preserves the
number of up spins, or equivalently the total $S^z$ charge, on the two lines.
\end{remark}

\section{Row-to-row transfer matrix}
\label{sec:six-vertex-row-transfer-matrix}

Consider a square lattice with $L$ columns and $M$ rows.  We impose periodic
boundary conditions in the horizontal direction.  A row configuration is specified
by the $L$ vertical arrows crossing a horizontal cut, hence the row Hilbert space
is
\begin{equation}
\label{eq:six-vertex-row-space}
    \calH_L=V_1\otimes V_2\otimes\cdots\otimes V_L,
    \qquad
    V_j\simeq\CC^2 .
\end{equation}
A row-to-row transfer matrix maps one row of vertical arrows to the next row,
summing over the internal horizontal arrows.

The efficient way to write this object is to introduce an auxiliary space
$V_a\simeq\CC^2$, which represents the horizontal arrow passing through the row.
Let $R_{aj}(u)$ act on the auxiliary space $a$ and the quantum space $j$.  The
monodromy matrix is
\begin{equation}
\label{eq:six-vertex-monodromy-matrix}
    M_a(u)
    =
    R_{aL}(u)R_{a,L-1}(u)\cdots R_{a1}(u),
\end{equation}
and the transfer matrix is the auxiliary trace
\begin{equation}
\label{eq:six-vertex-transfer-matrix}
    \tau(u)=\Tr_a M_a(u)
    =
    \Tr_a\!\bigl[R_{aL}(u)R_{a,L-1}(u)\cdots R_{a1}(u)\bigr] .
\end{equation}
The trace over $a$ closes the horizontal arrow into a periodic loop.  The
partition function on an $L\times M$ torus is then
\begin{equation}
\label{eq:six-vertex-partition-function-transfer}
    Z_{L,M}(u)=\Tr_{\calH_L}\tau(u)^M .
\end{equation}
Thus the thermodynamics of the two-dimensional classical model is controlled by
the spectrum of $\tau(u)$.

\begin{remark}[Two meanings of the same space]
In the classical model, $\calH_L$ is a convenient vector space for row
configurations.  In the quantum Hamiltonian limit, the same tensor product becomes
the Hilbert space of a length-$L$ spin chain.  This is the same conceptual move as
in the Ising transfer-matrix construction: a classical transfer direction becomes
imaginary time for a quantum problem.
\end{remark}

The ice rule implies a global $U(1)$ symmetry of the transfer matrix.  Let
\begin{equation}
\label{eq:six-vertex-total-magnetization}
    S^z_{\rm tot}=\frac12\sum_{j=1}^{L}\sigma_j^z .
\end{equation}
Since every local $R_{aj}(u)$ conserves the sum of the auxiliary and quantum
$z$-spins, tracing over the auxiliary space gives
\begin{equation}
\label{eq:six-vertex-transfer-U1-conservation}
    [\tau(u),S^z_{\rm tot}]=0 .
\end{equation}
The corresponding XXZ Hamiltonian therefore conserves total magnetization.

\section{The trigonometric six-vertex \texorpdfstring{$R$}{R}-matrix}
\label{sec:six-vertex-r-matrix}

A single transfer matrix is not yet enough to imply exact solvability.  The
integrable structure appears when the vertex weights are embedded in a
spectral-parameter family $R(u)$ satisfying the Yang--Baxter equation.  A standard
hyperbolic parametrization is
\begin{equation}
\label{eq:six-vertex-R-hyperbolic}
    R(u)=
    \begin{pmatrix}
        \sinh(u+\eta)&0&0&0\\
        0&\sinh u&\sinh\eta&0\\
        0&\sinh\eta&\sinh u&0\\
        0&0&0&\sinh(u+\eta)
    \end{pmatrix}.
\end{equation}
Here $u$ is the spectral parameter and $\eta$ is the anisotropy parameter.  In
this convention
\begin{equation}
\label{eq:six-vertex-hyperbolic-weights}
    a(u)=\sinh(u+\eta),
    \qquad
    b(u)=\sinh u,
    \qquad
    c(u)=\sinh\eta,
\end{equation}
and \cref{eq:six-vertex-Delta-from-abc} gives
\begin{equation}
\label{eq:xxz-Delta-cosh-eta}
    \Delta=\cosh\eta .
\end{equation}
For the critical regime one often uses the trigonometric parametrization
\begin{equation}
\label{eq:six-vertex-R-trigonometric}
    R(u)=
    \begin{pmatrix}
        \sin(u+\gamma)&0&0&0\\
        0&\sin u&\sin\gamma&0\\
        0&\sin\gamma&\sin u&0\\
        0&0&0&\sin(u+\gamma)
    \end{pmatrix},
    \qquad
    \Delta=\cos\gamma .
\end{equation}
The two forms are related by analytic continuation $\eta\leftrightarrow \ii\gamma$,
up to harmless scalar and phase normalizations.  The hyperbolic form is natural
for $\Delta>1$, while the trigonometric form is natural for $\abs{\Delta}<1$.

\begin{remark}[Scalar normalization]
The transformation
\begin{equation}
    R(u)\longmapsto f(u)R(u)
\end{equation}
with scalar $f(u)\neq0$ preserves the Yang--Baxter equation.  It multiplies the
transfer matrix by $f(u)^L$ and therefore changes only the scalar part of
$\log\tau(u)$.  In the Hamiltonian limit this produces an additive constant in
$H$, not a change of the interacting physics.
\end{remark}

The $R$-matrix in \cref{eq:six-vertex-R-hyperbolic} is \emph{regular}:
\begin{equation}
\label{eq:six-vertex-regularity}
    R(0)=\sinh\eta\,P,
\end{equation}
where $P:V\otimes V\to V\otimes V$ is the permutation operator
\begin{equation}
\label{eq:permutation-operator-definition}
    P(\ket{\alpha}\otimes\ket{\beta})
    =
    \ket{\beta}\otimes\ket{\alpha} .
\end{equation}
In Pauli matrices,
\begin{equation}
\label{eq:permutation-operator-pauli}
    P=\frac12\left(
        \id\otimes\id
        +\sigma^x\otimes\sigma^x
        +\sigma^y\otimes\sigma^y
        +\sigma^z\otimes\sigma^z
    \right).
\end{equation}
Regularity is the key property that makes the logarithmic derivative of
$\tau(u)$ local.

\section{Yang--Baxter equation and commuting transfer matrices}
\label{sec:six-vertex-yang-baxter}

Let $R_{ij}(u)$ denote $R(u)$ acting on the $i$th and $j$th tensor factors and as
identity on all others.  The Yang--Baxter equation is
\begin{equation}
\label{eq:six-vertex-YBE}
    R_{12}(u-v)R_{13}(u)R_{23}(v)
    =
    R_{23}(v)R_{13}(u)R_{12}(u-v) .
\end{equation}
It is an operator identity on $V_1\otimes V_2\otimes V_3$.  The six-vertex
$R$-matrix \cref{eq:six-vertex-R-hyperbolic} satisfies this equation.  In the
classical vertex-model language, \cref{eq:six-vertex-YBE} is a local
star--triangle move.  In the quantum scattering language, it is the consistency
condition for factorized three-body scattering.

\begin{theorem}[Commuting six-vertex transfer matrices]
\label{thm:six-vertex-commuting-transfer}
Let $R(u)$ satisfy the Yang--Baxter equation \cref{eq:six-vertex-YBE}.  Define
$\tau(u)$ by \cref{eq:six-vertex-transfer-matrix}.  Then
\begin{equation}
\label{eq:six-vertex-transfer-commutativity}
    [\tau(u),\tau(v)]=0,
    \qquad
    \text{for all }u,v .
\end{equation}
\end{theorem}

\begin{proof}
Introduce two auxiliary spaces $a$ and $b$.  The Yang--Baxter equation implies
the local relation
\begin{equation}
\label{eq:six-vertex-local-RLL}
    R_{ab}(u-v)R_{aj}(u)R_{bj}(v)
    =
    R_{bj}(v)R_{aj}(u)R_{ab}(u-v)
\end{equation}
for each quantum site $j$.  Multiplying this identity over all sites gives the
RTT relation
\begin{equation}
\label{eq:six-vertex-RTT}
    R_{ab}(u-v)M_a(u)M_b(v)
    =
    M_b(v)M_a(u)R_{ab}(u-v) .
\end{equation}
Taking $\Tr_a\Tr_b$ and using cyclicity of the trace in the auxiliary spaces,
$R_{ab}(u-v)$ can be moved through the trace and canceled.  The result is
\begin{equation}
    \Tr_a M_a(u)\,\Tr_b M_b(v)
    =
    \Tr_b M_b(v)\,\Tr_a M_a(u),
\end{equation}
which is exactly $\tau(u)\tau(v)=\tau(v)\tau(u)$.
\end{proof}

The implication is substantial.  Expanding the logarithm of $\tau(u)$ around a
regular point gives an infinite sequence of mutually commuting conserved charges:
\begin{equation}
\label{eq:six-vertex-log-transfer-charges}
    \log\tau(u)
    =
    \sum_{n=0}^{\infty}(u-u_0)^n Q_n,
    \qquad
    [Q_m,Q_n]=0 .
\end{equation}
The first nontrivial local charge is the XXZ Hamiltonian.  Higher charges have
larger spatial range and are responsible for integrability.

\begin{principle}[Yang--Baxter integrability]
The Yang--Baxter equation is the local algebraic identity that upgrades one
transfer matrix into a commuting one-parameter family.  The logarithm of this
family generates the Hamiltonian and the higher conserved charges.
\end{principle}

\section{Hamiltonian limit and the XXZ chain}
\label{sec:six-vertex-hamiltonian-limit}

We now derive the quantum Hamiltonian associated with the regular six-vertex
transfer matrix.  For a homogeneous periodic chain, regularity gives
\begin{equation}
\label{eq:six-vertex-transfer-at-zero-shift}
    \tau(0)
    =
    (\sinh\eta)^L\,U,
\end{equation}
where $U$ is the one-site translation operator on the spin chain.  Taking the
logarithmic derivative removes this translation and produces a nearest-neighbor
operator:
\begin{equation}
\label{eq:six-vertex-log-derivative-locality}
    \left.\frac{\dd}{\dd u}\log\tau(u)\right|_{u=0}
    =
    \sum_{j=1}^{L}
    \frac{1}{\sinh\eta}
    P_{j,j+1}R'_{j,j+1}(0),
\end{equation}
with $j+1$ understood modulo $L$.  The local operator can be computed explicitly.
From \cref{eq:six-vertex-R-hyperbolic},
\begin{equation}
\label{eq:six-vertex-R-prime-zero}
    R'(0)=
    \begin{pmatrix}
        \cosh\eta&0&0&0\\
        0&1&0&0\\
        0&0&1&0\\
        0&0&0&\cosh\eta
    \end{pmatrix}.
\end{equation}
Therefore
\begin{equation}
\label{eq:six-vertex-local-log-derivative-pauli}
    \frac{1}{\sinh\eta}P R'(0)
    =
    \frac{1}{2\sinh\eta}
    \left(
        \sigma^x\otimes\sigma^x
        +\sigma^y\otimes\sigma^y
        +\cosh\eta\,\sigma^z\otimes\sigma^z
        +\cosh\eta\,\id\otimes\id
    \right).
\end{equation}
Multiplying the logarithmic derivative by $2\sinh\eta$ and subtracting an
additive constant gives
\begin{equation}
\label{eq:xxz-from-six-vertex-integrability-normalization}
    H_{\rm XXZ}^{(+)}
    =
    J\sum_{j=1}^{L}
    \left(
        \sigma_j^x\sigma_{j+1}^x
        +\sigma_j^y\sigma_{j+1}^y
        +\Delta\bigl(\sigma_j^z\sigma_{j+1}^z-1\bigr)
    \right),
    \qquad
    \Delta=\cosh\eta .
\end{equation}
Equivalently,
\begin{equation}
\label{eq:xxz-hamiltonian-log-derivative-final}
    H_{\rm XXZ}^{(+)}
    =
    2J\sinh\eta
    \left.\frac{\dd}{\dd u}\log\tau(u)\right|_{u=0}
    -2J\Delta L .
\end{equation}
The superscript $(+)$ indicates the standard antiferromagnetic integrability
normalization for $J>0$ and $\Delta>0$.  The constant $-2J\Delta L$ fixes the
energy of the fully polarized reference state to zero; it has no effect on
eigenvectors, gaps, or correlation functions.

In spin-operator normalization, $S_j^\alpha=\sigma_j^\alpha/2$, the same
Hamiltonian is
\begin{equation}
\label{eq:xxz-spin-operator-normalization-from-six-vertex}
    H_{\rm XXZ}^{(+)}
    =
    4J\sum_{j=1}^{L}
    \left(
        S_j^xS_{j+1}^x
        +S_j^yS_{j+1}^y
        +\Delta S_j^zS_{j+1}^z
    \right)
    -J\Delta L .
\end{equation}
Only the overall energy scale and additive constant differ from other common XXZ
conventions.

\begin{remark}[Relation to the earlier Pauli convention]
Earlier chapters used the ferromagnetic Pauli convention
\begin{equation}
    H_{\rm XXZ}^{\rm ferro}
    =
    -\frac{J_{\rm f}}{2}\sum_j
    \left(
        \sigma_j^x\sigma_{j+1}^x
        +\sigma_j^y\sigma_{j+1}^y
        +\Delta\sigma_j^z\sigma_{j+1}^z
    \right)
    -h\sum_j\sigma_j^z .
\end{equation}
The Hamiltonian generated directly by the six-vertex $R$-matrix is instead the
integrability convention \cref{eq:xxz-from-six-vertex-integrability-normalization}.
The two are related by an overall sign and energy-scale choice, plus an additive
constant.  Whenever Bethe-ansatz formulas are quoted, the sign convention must be
stated explicitly.
\end{remark}

The trigonometric $R$-matrix gives the same formula with
\begin{equation}
\label{eq:xxz-gapless-Delta-cos-gamma}
    \Delta=\cos\gamma,
    \qquad
    \sinh\eta\longrightarrow\sin\gamma .
\end{equation}
Thus the critical XXZ regime $\abs{\Delta}<1$ is naturally parametrized by
$0<\gamma<\pi$.

\section{Why the Hamiltonian is local}
\label{sec:six-vertex-why-local}

It is useful to isolate the mechanism behind locality.  The transfer matrix
\begin{equation}
    \tau(u)=\Tr_a R_{aL}(u)\cdots R_{a1}(u)
\end{equation}
is a nonlocal operator on the whole chain.  Nevertheless, at the regular point
$u=0$ every $R_{aj}(0)$ is a permutation.  The product of permutations simply
moves the auxiliary spin through the chain and becomes the translation operator
after the auxiliary trace.  Differentiating with respect to $u$ replaces exactly
one local permutation by $R'(0)$, while all other factors remain permutations.
After multiplying by $\tau(0)^{-1}$, this single defect becomes a local
nearest-neighbor density.

This is the vertex-model analogue of the Trotter logic in the Ising--TFIM
correspondence.  The difference is that the present transfer matrix is not merely
one short-time evolution operator.  It belongs to a commuting spectral-parameter
family, and this additional structure produces infinitely many conserved charges.

\begin{proposition}[Local charge from regularity]
\label{prop:regular-transfer-local-charge}
Let $R(u)$ satisfy $R(0)=\rho P$ with $\rho\neq0$, and define the homogeneous
periodic transfer matrix by \cref{eq:six-vertex-transfer-matrix}.  Then
\begin{equation}
    \left.\frac{\dd}{\dd u}\log\tau(u)\right|_{u=0}
    =
    \sum_{j=1}^{L}h_{j,j+1},
    \qquad
    h_{j,j+1}=\rho^{-1}P_{j,j+1}R'_{j,j+1}(0),
\end{equation}
up to the conventional orientation of the one-site translation operator.
\end{proposition}

\section{Conserved magnetization and interacting character}
\label{sec:six-vertex-conserved-magnetization-interaction}

The XXZ Hamiltonian obtained above satisfies
\begin{equation}
\label{eq:xxz-U1-symmetry}
    [H_{\rm XXZ}^{(+)},S^z_{\rm tot}]=0 .
\end{equation}
Using $\sigma_j^\pm=(\sigma_j^x\pm\ii\sigma_j^y)/2$, one may rewrite the exchange
part as
\begin{equation}
\label{eq:xxz-ladder-form}
    \sigma_j^x\sigma_{j+1}^x+
    \sigma_j^y\sigma_{j+1}^y
    =
    2\left(
        \sigma_j^+\sigma_{j+1}^-
        +
        \sigma_j^-\sigma_{j+1}^+
    \right).
\end{equation}
Thus the $XY$ part hops a down spin through an up-spin background, while the
$\Delta\sigma_j^z\sigma_{j+1}^z$ term gives a nearest-neighbor interaction between
magnons.  At $\Delta=0$ the model is the XX chain and becomes free after the
Jordan--Wigner transformation.  At generic $\Delta$, the Jordan--Wigner image
contains a density--density interaction and is not a quadratic fermion model.

Explicitly, under Jordan--Wigner fermionization, the $XY$ part becomes hopping,
whereas
\begin{equation}
\label{eq:xxz-density-density-interaction}
    \sigma_j^z\sigma_{j+1}^z
    =
    (1-2n_j)(1-2n_{j+1})
    =
    1-2(n_j+n_{j+1})+4n_jn_{j+1} .
\end{equation}
The quartic term $n_jn_{j+1}$ is the interacting part.  Therefore the XXZ chain
is a genuinely interacting integrable model, except at the free-fermion point
$\Delta=0$.

\begin{remark}[Why this chapter leaves the spectrum unfinished]
For the Ising, XY, Kitaev, and SSH chains, exact solvability meant reducing the
Hamiltonian to independent BdG blocks.  For the XXZ chain this is impossible at
generic $\Delta$ because of the quartic Jordan--Wigner interaction.  Its solution
requires Bethe ansatz: many-body eigenstates are built from interacting magnons
whose scattering factorizes.  The six-vertex transfer matrix supplies the
Yang--Baxter reason why that factorization is consistent.
\end{remark}

\section{Anisotropy regimes}
\label{sec:six-vertex-xxz-regimes}

For the antiferromagnetic XXZ convention
\begin{equation}
\label{eq:xxz-standard-spin-regime-convention}
    H_{\rm XXZ}^{S}
    =
    \mathcal{J}\sum_j
    \left(
        S_j^xS_{j+1}^x
        +S_j^yS_{j+1}^y
        +\Delta S_j^zS_{j+1}^z
    \right),
    \qquad
    \mathcal{J}>0,
\end{equation}
the familiar zero-field phase structure is summarized as follows:
\begin{center}
\begin{tabular}{@{}L{0.22\linewidth}L{0.30\linewidth}L{0.34\linewidth}@{}}
\toprule
Parameter regime & Useful parametrization & Quantum-chain behavior \\
\midrule
$\Delta>1$ & $\Delta=\cosh\eta$ & gapped antiferromagnetic, or N\'{e}el, regime \\
$-1<\Delta<1$ & $\Delta=\cos\gamma$ & gapless critical Luttinger-liquid regime \\
$\Delta=1$ & isotropic XXX point & critical with full $SU(2)$ symmetry and marginal corrections \\
$\Delta<-1$ & analytic continuation or sign-twisted convention & gapped ferromagnetic regime \\
\bottomrule
\end{tabular}
\end{center}
For the classical six-vertex model with positive weights, one also classifies
regions using
\begin{equation}
    \Delta=\frac{a^2+b^2-c^2}{2ab} .
\end{equation}
The regime $\abs{\Delta}<1$ is the critical disordered phase of the classical
vertex model.  The regions $\Delta>1$ and $\Delta<-1$ correspond to ordered
ferroelectric and antiferroelectric regimes, respectively, depending on which
weights dominate.  The exact terminology is classical-model dependent, but the
same invariant $\Delta$ controls the integrable family.

\begin{remark}[Classical phases versus quantum phases]
The two-dimensional six-vertex partition function at a fixed spectral parameter
and the one-dimensional XXZ Hamiltonian obtained by a logarithmic derivative are
not the same physical object.  They are two different uses of the same commuting
transfer-matrix family.  The shared parameter $\Delta$ explains why their phase
diagrams are closely related, but one should not identify a classical ordered
region with a quantum ground-state phase without specifying the Hamiltonian limit
and the sign convention.
\end{remark}

\section{What the transfer matrix solution will become}
\label{sec:six-vertex-what-comes-next}

The six-vertex transfer matrix is diagonalized by Bethe ansatz.  In the algebraic
form, one writes the monodromy matrix as a $2\times2$ matrix in auxiliary space,
\begin{equation}
\label{eq:six-vertex-ABA-monodromy-entries}
    M_a(u)=
    \begin{pmatrix}
        A(u)&B(u)\\
        C(u)&D(u)
    \end{pmatrix}_a,
    \qquad
    \tau(u)=A(u)+D(u).
\end{equation}
The operator $B(u)$ creates Bethe quasiparticles above a reference state, and the
RTT relation generated by the Yang--Baxter equation gives the algebra needed to
move $A(u)$ and $D(u)$ through a product of $B$'s.  Requiring unwanted terms to
vanish gives the Bethe equations.

In the coordinate form, the same result appears as factorized magnon scattering.
In a sector with $M$ down spins, a wave function is a sum over plane waves in each
ordered coordinate region,
\begin{equation}
    \Psi(x_1,\ldots,x_M)
    =
    \sum_{P\in S_M}A(P)
    \exp\left(\ii\sum_{a=1}^{M}k_{P_a}x_a\right),
    \qquad
    x_1<\cdots<x_M .
\end{equation}
The two-body scattering phase determines how the amplitudes $A(P)$ change when
neighboring momenta are exchanged.  Periodic boundary conditions then lead to
Bethe equations.  The following chapters derive these statements explicitly.

\section{Summary}
\label{sec:six-vertex-summary}

The essential chain of ideas is
\begin{equation}
\label{eq:six-vertex-summary-chain}
\boxed{
\begin{gathered}
    \text{ice-rule vertex weights}
    \;\Longrightarrow\;
    R(u)\text{ with }U(1)\text{ conservation}
    \\
    \Longrightarrow\;
    \text{Yang--Baxter equation}
    \;\Longrightarrow\;
    [\tau(u),\tau(v)]=0
    \\
    \Longrightarrow\;
    H_{\rm XXZ}\propto
    \left.\frac{\dd}{\dd u}\log\tau(u)\right|_{u=0}
    \;\Longrightarrow\;
    \text{interacting Bethe-ansatz integrability.}
\end{gathered}}
\end{equation}
The six-vertex model is therefore the canonical bridge between two-dimensional
classical vertex models and one-dimensional interacting quantum spin chains.
Unlike the Ising--TFIM correspondence, the resulting quantum chain is not solved
by free fermions at generic anisotropy.  Its solvability is instead controlled by
Yang--Baxter integrability and factorized scattering.
